\chapter{ADR-027 --- Kiến trúc tầng module \& Mobile core (2026-09-17)}
% Hình của ADR-027 — \input từ 74-adr027-module-layering.tex
\tikzset{
  a27box/.style={draw=black!60, rounded corners=2pt, font=\ttfamily\scriptsize, inner sep=3pt, align=center, fill=white},
  a27new/.style={a27box, dashed, draw=blue!70!black, fill=blue!7},
  a27mv/.style={a27box, dashed, draw=orange!80!black, fill=orange!12},
  a27bad/.style={a27box, draw=red!70!black, fill=red!8},
  a27dep/.style={-{Stealth[length=5pt]}, thick, draw=black!55},
  a27badep/.style={-{Stealth[length=5pt]}, thick, draw=red!70!black},
  a27band/.style={rounded corners=4pt, draw=none, fill=#1},
  a27lbl/.style={anchor=east, align=right, font=\sffamily\scriptsize\bfseries},
  a27mdl/.style={draw=black!55, fill=white, rounded corners=1pt, font=\ttfamily\tiny, inner sep=2pt, minimum width=2.9cm},
  a27hub/.style={a27mdl, fill=yellow!30, very thick, font=\ttfamily\scriptsize\bfseries},
  a27ext/.style={a27mdl, dashed, fill=black!3},
  a27grp/.style={draw=#1!70!black, fill=#1!6, rounded corners=3pt, inner sep=5pt},
  a27gl/.style={anchor=south west, font=\sffamily\tiny\bfseries, inner sep=1pt},
  a27rel/.style={-{Stealth[length=3.5pt]}, draw=black!55},
  a27xrel/.style={-{Stealth[length=5pt]}, draw=orange!80!black, thick},
}

\newcommand{\adrLayerFigure}{%
\begin{figure}[!htbp]
\centering
\resizebox{\textwidth}{!}{%
\begin{tikzpicture}[x=1cm,y=1cm]
  \fill[a27band=black!6]  (0,0)   rectangle (19.5,1.0);
  \fill[a27band=green!7]  (0,1.4) rectangle (19.5,2.2);
  \fill[a27band=green!12] (0,2.6) rectangle (12.4,6.0);
  \fill[a27band=blue!6]   (12.8,2.6) rectangle (19.5,6.0);
  \fill[a27band=blue!10]  (0,6.4) rectangle (19.5,9.6);
  \node[a27lbl] at (-0.2,0.5)  {L0\\Odoo + bên thứ 3};
  \node[a27lbl] at (-0.2,1.8)  {L1\\Nền};
  \node[a27lbl] at (-0.2,4.3)  {L2\\Nghiệp vụ | khung kênh\\(cùng cấp)};
  \node[a27lbl] at (-0.2,8.0)  {L3\\Module ghép\\(a: nền · b: phân hệ)};

  \node[a27box, minimum width=17cm] at (9.75,0.5)
    {Odoo 19 (base · web · mail · portal · sale · stock · account) \ \ | \ \ muk\_web\_theme \ \ | \ \ bên thứ 3};
  \node[a27box] (core) at (2.5,1.8) {wujia\_core {\sffamily\tiny area, ward}};
  \node[a27new] (uicore) at (6.6,1.8) {ui\_core {\sffamily\tiny chỉ tạo khi 2 kênh cần}};
  \node[a27new] (plat) at (11.0,1.8) {*\_core {\sffamily\tiny security · approval · job · audit}};

  \foreach \x/\n in {0.9/franchise, 2.45/\_contract, 4.0/\_inspection, 5.55/\_operations,
                     7.1/sale, 8.65/fleet, 10.2/delivery, 11.75/account}
    \node[a27box, font=\ttfamily\tiny, minimum height=0.6cm, inner sep=2pt] at (\x,5.0) {\n};
  \node[font=\sffamily\tiny, anchor=north west] at (0.2,5.95) {nghiệp vụ đang sạch};
  \foreach \x/\n in {0.95/order\\\_window, 2.75/info\\\_request, 4.55/knowledge, 6.35/support, 8.15/announce-\\ment, 9.95/exam, 11.65/return}
    \node[a27mv, minimum height=0.8cm] at (\x,3.4) {\n};
  \node[font=\sffamily\tiny, text=orange!80!black, anchor=south west] at (0.2,2.62) {tách ra khỏi wujia\_portal\_<x>};

  \node[a27box] (play) at (14.5,4.3) {portal\_layout\\{\sffamily\tiny Vuexy · assets\_frontend}};
  \node[a27new] (mcore) at (17.7,4.3) {mobile\_core\\{\sffamily\tiny MuK · assets\_backend}};
  \node[font=\sffamily\tiny\bfseries, anchor=north west] at (13.0,5.95) {khung kênh (không biết nghiệp vụ)};
  \draw[black!50, dashed, thick] (12.6,2.7) -- (12.6,5.9);

  \node[a27box] (pbase) at (9.0,7.2) {portal\_base\\{\sffamily\tiny nền ghép portal: store · Home · bus}};
  \node[font=\sffamily\tiny\bfseries, anchor=east, text=blue!50!black] at (pbase.west) {L3a\ };
  \node[font=\sffamily\tiny\bfseries, anchor=south west, text=blue!50!black] at (0.2,9.1) {L3b};
  \node[a27box, minimum width=6cm] (px) at (4.2,8.8)
    {portal\_<x> \quad controller + QWeb\\{\sffamily\tiny sale · debt · delivery · exam · return · \ldots}};
  \node[a27new, minimum width=4.2cm] (mx) at (17.2,8.8)
    {mobile\_<x> \quad auto\_install\\{\sffamily\tiny chỉ tạo khi có view riêng}};

  \draw[a27dep] (px.east) -| (pbase.north);
  \draw[a27dep] (2.0,8.25) -- (2.0,6.02);
  \draw[a27dep] (pbase.south) -- (9.0,6.02);
  \draw[a27dep] (pbase.east) -| (play.north);
  \draw[a27dep] (mx.south -| mcore) -- (mcore.north);
  \draw[a27dep] (mx.west) -| (11.9,6.02);
  \draw[a27dep] (play.south) -- (14.5,1.02);
  \draw[a27dep] (mcore.south) -- (17.7,1.02);
  \draw[a27dep] (2.5,2.6) -- (core.north);
  \draw[a27dep] (core.south) -- (2.5,1.02);

  \node[a27box, anchor=west] at (0,-0.7) {hiện có};
  \node[a27mv, anchor=west] at (1.6,-0.7) {tách ra};
  \node[a27new, anchor=west] at (3.1,-0.7) {tạo mới};
  \node[font=\sffamily\scriptsize, anchor=west] at (4.8,-0.7)
    {Mũi tên $=$ \texttt{depends}, chỉ đi xuống. Bỏ tiền tố \texttt{wujia\_}. Nghiệp vụ và khung kênh không depend nhau; portal\_* và mobile\_* không depend chéo.};
\end{tikzpicture}}
\caption{ADR-027 --- bốn tầng module: nghiệp vụ và khung kênh cùng cấp, chỉ module ghép mới thấy cả hai.}
\label{fig:adr27-layers}
\end{figure}}

\newcommand{\adrWhyFigure}{%
\begin{figure}[htbp]
\centering
\resizebox{0.92\textwidth}{!}{%
\begin{tikzpicture}[x=1cm,y=1cm, node distance=0.55cm]
  \node[font=\sffamily\small\bfseries] at (2.2,1.0) {Hiện tại};
  \node[a27new] (m1) at (2.2,0) {wujia\_mobile\_exam};
  \node[a27bad, below=of m1] (p1) {wujia\_portal\_exam\\{\sffamily\tiny model thi + view backend + QWeb}};
  \node[a27bad, below=of p1] (b1) {wujia\_portal\_base};
  \node[a27bad, below=of b1] (l1) {wujia\_portal\_layout\\{\sffamily\tiny web.assets\_frontend}};
  \node[a27box, right=0.5cm of b1] (s1) {wujia\_sale};
  \draw[a27badep] (m1) -- (p1);
  \draw[a27badep] (p1) -- (b1);
  \draw[a27badep] (b1) -- (l1);
  \draw[a27badep] (b1) -- (s1);
  \node[font=\sffamily\scriptsize, text=red!70!black, align=center, below=0.3cm of l1]
    {mobile kéo theo cả portal\\không tách repo được};

  \node[font=\sffamily\small\bfseries] at (11.9,1.0) {Sau khi tách};
  \node[a27new] (m2) at (10.1,0) {wujia\_mobile\_exam};
  \node[a27box] (p2) at (13.7,0) {wujia\_portal\_exam\\{\sffamily\tiny controller + QWeb}};
  \node[a27mv] (d2) at (11.9,-2.2) {wujia\_exam\\{\sffamily\tiny model + view backend + quyền}};
  \node[a27new] (c2) at (8.4,-2.2) {wujia\_mobile\_core};
  \node[a27box] (b2) at (15.4,-2.2) {wujia\_portal\_base};
  \node[a27box] (f2) at (11.9,-4.2) {wujia\_franchise};
  \draw[a27dep] (m2) -- (d2);
  \draw[a27dep] (m2) -- (c2);
  \draw[a27dep] (p2) -- (d2);
  \draw[a27dep] (p2) -- (b2);
  \draw[a27dep] (d2) -- (f2);
  \node[font=\sffamily\scriptsize, text=green!40!black, align=center] at (11.9,-5.3)
    {mobile và portal chỉ gặp nhau ở tầng nghiệp vụ};
\end{tikzpicture}}
\caption{Vì sao phải tách nghiệp vụ khỏi module portal --- ví dụ phân hệ Đăng ký thi.}
\label{fig:adr27-why}
\end{figure}}

\newcommand{\adrErdFigure}{%
\begin{figure}[!htbp]
\centering
\resizebox{!}{0.92\textheight}{%
\begin{tikzpicture}[x=1cm,y=1cm, node distance=2pt]
  % ---- wujia_franchise (hub)
  \node[a27mdl] (fdoc) at (0,1.0) {franchise.document};
  \node[a27hub, below=6pt of fdoc] (hub) {franchise.management};
  \node[a27mdl, below=6pt of hub] (fmem) {franchise.member};
  \begin{scope}[on background layer]\node[a27grp=green, fit=(fdoc)(hub)(fmem)] (Gfr) {};\end{scope}
  \node[a27gl] at (Gfr.north west) {wujia\_franchise};
  \draw[a27rel] (fdoc) -- (hub);
  \draw[a27rel] (fmem) -- (hub);

  % ---- wujia_core
  \node[a27mdl] (area) at (0,6.0) {res.area};
  \node[a27mdl, below=of area] (ward) {res.ward};
  \begin{scope}[on background layer]\node[a27grp=green, fit=(area)(ward)] (Gco) {};\end{scope}
  \node[a27gl] at (Gco.north west) {wujia\_core};

  % ---- contract
  \node[a27mdl] (ctr) at (-5.2,6.0) {franchise.contract};
  \begin{scope}[on background layer]\node[a27grp=green, fit=(ctr)] (Gct) {};\end{scope}
  \node[a27gl] at (Gct.north west) {wujia\_franchise\_contract};

  % ---- account
  \node[a27ext] (mv) at (5.2,6.0) {account.move*};
  \node[a27ext, below=of mv] (pay) {account.payment*};
  \begin{scope}[on background layer]\node[a27grp=green, fit=(mv)(pay)] (Gac) {};\end{scope}
  \node[a27gl] at (Gac.north west) {wujia\_account};

  % ---- operations
  \node[a27mdl] (emp) at (-5.2,2.2) {franchise.employee};
  \node[a27mdl, below=of emp] (asg) {employee.assignment};
  \node[a27mdl, below=of asg] (wsch) {work.schedule};
  \node[a27mdl, below=of wsch] (shift) {shift.template};
  \node[a27mdl, below=of shift] (rev) {franchise.revenue};
  \node[a27mdl, below=of rev] (exp) {franchise.expense};
  \begin{scope}[on background layer]\node[a27grp=green, fit=(emp)(exp)] (Gop) {};\end{scope}
  \node[a27gl] at (Gop.north west) {wujia\_franchise\_operations};
  \draw[a27rel] (asg.west) to[bend left=40] (emp.west);
  \draw[a27rel] (wsch.west) to[bend left=40] (asg.west);

  % ---- inspection
  \node[a27mdl] (ins) at (5.2,2.2) {franchise.inspection};
  \node[a27mdl, below=of ins] (insl) {inspection.line};
  \node[a27mdl, below=of insl] (tpl) {inspection.template};
  \node[a27mdl, below=of tpl] (grd) {inspection.grade};
  \node[a27mdl, below=of grd] (sch) {supervision.schedule};
  \begin{scope}[on background layer]\node[a27grp=green, fit=(ins)(sch)] (Gin) {};\end{scope}
  \node[a27gl] at (Gin.north west) {wujia\_franchise\_inspection};
  \draw[a27rel] (insl.east) to[bend right=40] (ins.east);
  \draw[a27rel] (sch.east) to[bend right=40] (grd.east);

  % ---- sale / delivery / fleet
  \node[a27ext] (so) at (-5.2,-4.2) {sale.order*};
  \node[a27mdl, below=of so] (pcat) {product.category};
  \node[a27mdl, below=of pcat] (sdr) {supply.demand.report};
  \begin{scope}[on background layer]\node[a27grp=green, fit=(so)(sdr)] (Gsa) {};\end{scope}
  \node[a27gl] at (Gsa.north west) {wujia\_sale};

  \node[a27ext] (batch) at (0,-4.2) {stock.picking.batch*};
  \node[a27ext, below=of batch] (pick) {stock.picking*};
  \begin{scope}[on background layer]\node[a27grp=green, fit=(batch)(pick)] (Gde) {};\end{scope}
  \node[a27gl] at (Gde.north west) {wujia\_delivery};

  \node[a27mdl] (fl) at (5.2,-4.2) {fleet.management};
  \node[a27mdl, below=of fl] (prov) {fleet.provider};
  \node[a27mdl, below=of prov] (ftype) {fleet.type};
  \node[a27mdl, below=of ftype] (fpl) {fleet.pricelist};
  \begin{scope}[on background layer]\node[a27grp=green, fit=(fl)(fpl)] (Gfl) {};\end{scope}
  \node[a27gl] at (Gfl.north west) {wujia\_fleet};
  \draw[a27rel] (batch) -- (fl);
  \draw[a27rel] (so) -- (batch);

  % ---- phân hệ tách ra (L2 mới)
  \node[a27mdl] (ow) at (-5.2,-8.6) {order.window};
  \begin{scope}[on background layer]\node[a27grp=orange, fit=(ow)] (Gow) {};\end{scope}
  \node[a27gl] at (Gow.north west) {wujia\_order\_window};

  \node[a27mdl] (ir) at (0,-8.6) {info.update.request};
  \begin{scope}[on background layer]\node[a27grp=orange, fit=(ir)] (Gir) {};\end{scope}
  \node[a27gl] at (Gir.north west) {wujia\_info\_request};

  \node[a27mdl] (art) at (5.2,-8.6) {knowledge.article};
  \node[a27mdl, below=of art] (kcat) {knowledge.category};
  \node[a27mdl, below=of kcat] (ktag) {knowledge.tag};
  \begin{scope}[on background layer]\node[a27grp=orange, fit=(art)(ktag)] (Gkn) {};\end{scope}
  \node[a27gl] at (Gkn.north west) {wujia\_knowledge};
  \draw[a27rel] (art.east) to[bend right=40] (kcat.east);

  \node[a27mdl] (tk) at (-5.2,-11.4) {support.ticket};
  \node[a27mdl, below=of tk] (tcat) {support.category};
  \begin{scope}[on background layer]\node[a27grp=orange, fit=(tk)(tcat)] (Gsu) {};\end{scope}
  \node[a27gl] at (Gsu.north west) {wujia\_support};

  \node[a27mdl] (nt) at (0,-11.4) {notification};
  \node[a27mdl, below=of nt] (nrd) {notification.read};
  \node[a27mdl, below=of nrd] (ntyp) {notification.type};
  \begin{scope}[on background layer]\node[a27grp=orange, fit=(nt)(ntyp)] (Gno) {};\end{scope}
  \node[a27gl] at (Gno.north west) {wujia\_notification};
  \draw[a27rel] (nrd.west) to[bend left=40] (nt.west);

  \node[a27mdl] (reg) at (5.2,-11.4) {exam.registration};
  \node[a27mdl, below=of reg] (regl) {registration.line};
  \node[a27mdl, below=of regl] (ses) {exam.session};
  \node[a27mdl, below=of ses] (crs) {exam.course};
  \node[a27mdl, below=of crs] (slot) {exam.time.slot};
  \begin{scope}[on background layer]\node[a27grp=orange, fit=(reg)(slot)] (Gex) {};\end{scope}
  \node[a27gl] at (Gex.north west) {wujia\_exam};
  \draw[a27rel] (regl.east) to[bend right=40] (reg.east);
  \draw[a27rel] (ses.east) to[bend right=40] (crs.east);

  \node[a27mdl] (rr) at (-5.2,-15.0) {return.request};
  \node[a27mdl, below=of rr] (alloc) {compensation.allocation};
  \node[a27mdl, below=of alloc] (itype) {return.issue.type};
  \begin{scope}[on background layer]\node[a27grp=orange, fit=(rr)(itype)] (Gre) {};\end{scope}
  \node[a27gl] at (Gre.north west) {wujia\_return};
  \draw[a27rel] (alloc.west) to[bend left=40] (rr.west);

  % ---- model riêng của kênh portal (L4)
  \node[a27mdl] (cart) at (5.2,-18.2) {portal.cart (+line)};
  \node[a27mdl, below=of cart] (debt) {portal.debt (abstract)};
  \begin{scope}[on background layer]\node[a27grp=blue, fit=(cart)(debt)] (Gpo) {};\end{scope}
  \node[a27gl] at (Gpo.north west) {L4 · model riêng của portal};

  % ---- quan hệ xuyên module (franchise_id / area_id / sale_order_id)
  \draw[a27rel] (Gfr.north) -- (Gco.south);
  \draw[a27rel] (Gct.south east) -- (hub.north west);
  \draw[a27rel] (Gac.south west) -- (hub.north east);
  \draw[a27rel] (Gop.east) -- (hub.west);
  \draw[a27rel] (Gin.west) -- (hub.east);
  \draw[a27rel] (Gsa.north) -- node[font=\sffamily\tiny, fill=white, inner sep=1pt, pos=0.55]{franchise\_id} (Gfr.south west);
  \draw[a27rel] (Gde.north) -- node[font=\sffamily\tiny, fill=white, inner sep=1pt]{franchise\_id} (Gfr.south);
  \node[draw=orange!80!black, dashed, thick, rounded corners=5pt, inner sep=9pt, fit=(Gow)(Gkn)(Gre)(Gex)] (Gmv) {};
  \draw[a27xrel] (Gmv.north -| -2.9,0) -- (-2.9,0.2) -- (Gfr.west |- 0,0.2)
     node[pos=0.4, font=\sffamily\tiny, fill=white, inner sep=1pt, rotate=90]{franchise\_id (trừ knowledge)};
  \draw[a27xrel] (Gmv.north -| -6.3,0) -- (Gsa.south -| -6.3,0)
     node[midway, font=\sffamily\tiny, fill=white, inner sep=1pt, rotate=90]{sale\_order\_id};
  \node[font=\sffamily\scriptsize, align=left, anchor=north west, text width=8.8cm] at (-7.2,-17.4)
    {\textbf{Đọc hình.} Hộp xanh lá $=$ module nghiệp vụ đang sạch; cam $=$ phân hệ tách ra khỏi
     \texttt{wujia\_portal\_<x>}; xanh dương $=$ model riêng của kênh. Nền vàng là \emph{hub}
     \texttt{wujia.franchise.management}. Model có \texttt{*} là model Odoo được kế thừa. Mũi tên
     đen $=$ Many2one. Khung cam đứt nét gom 7 phân hệ tách ra; hai mũi tên cam là \textbf{quan hệ gộp}: mọi phân
     hệ trong khung có \texttt{franchise\_id} $\rightarrow$ hub (trừ knowledge); \texttt{support},
     \texttt{return} trỏ \texttt{sale.order} và \texttt{stock.picking.batch};
     \texttt{order\_window} trỏ \texttt{res.area}; \texttt{exam}, \texttt{notification.read}
     trỏ thêm \texttt{franchise.member}.};
\end{tikzpicture}}
\caption{Phác thảo ERD theo module đích --- model chính và quan hệ Many2one then chốt (đo từ source \texttt{12750a2}).}
\label{fig:adr27-erd}
\end{figure}}

\sloppy

\emph{Chapter này là một \textbf{quyết định kiến trúc đề xuất}, chưa có dòng code nghiệp vụ nào.
Nó trả lời câu hỏi của BA và chủ dự án: muốn dựng \emph{mobile core} cho backend ERP thì
module phải phụ thuộc nhau thế nào, để sau này tách Wujia thành các folder/repo riêng như
beesuite mà không bị kéo ngược lên portal. Mọi con số trong chapter đo trên \texttt{main}
\texttt{12750a2} (sau khi pull 109 commit) và trên UAT chỉ-đọc qua XML-RPC.}

\section{Bối cảnh --- vì sao phải chốt trước khi viết module mobile đầu tiên}

Khi mở backend ERP trên điện thoại (ví dụ view kanban), thao tác vẫn là kiểu máy tính: dựa
vào \emph{hover}, không có trạng thái chạm, không có chọn bản ghi như một app mobile. BA đề
xuất dựng \textbf{một mobile core dùng chung}: chỉ sửa CSS/JS của backend, không đụng portal,
không đụng controller.

Chủ dự án nêu ra đúng chỗ đáng lo. Nếu mỗi phân hệ có một module mobile
(\texttt{wujia\_mobile\_sale}, \texttt{wujia\_mobile\_franchise}\ldots) thì module đó phải depend
cả mobile core lẫn phân hệ nghiệp vụ. Nếu phân hệ nghiệp vụ lại đang depend portal, thì
\emph{mobile bị kéo ngược lên portal}. Hai kênh dính vào nhau, và không tách repo được nữa.

Quyết định dependency là loại quyết định \textbf{rất đắt để sửa về sau}: đổi chủ một model giữa
hai module trên DB đang chạy phải viết migration \texttt{ir\_model\_data} (xem \S\ref{sec:adr27-xmlid}).
Vì vậy phải chốt trước.

\section{Hiện trạng --- bản đồ phụ thuộc thật}

\subsection*{Tên gọi dễ nhầm}

\begin{itemize}
\item \textbf{Không có module \texttt{wujia\_base}.} \texttt{base} là module lõi của Odoo. Cái
  hay được gọi tắt là ``base'' là \texttt{wujia\_portal\_base}, tức lõi của portal.
\item \texttt{wujia\_core} rất mỏng: chỉ 2 model master data \texttt{res.area} (khu vực) và
  \texttt{res.ward} (phường), depend \texttt{base, mail, contacts}.
\item \texttt{wujia\_franchise} depend module \texttt{portal} \emph{của Odoo}. Đó là để có nhóm
  \texttt{base.group\_portal} viết ACL/rule, \textbf{không phải} depend portal Vuexy của Wujia.
\end{itemize}

\subsection*{Xếp 26 module Wujia vào tầng (đo từ \texttt{\_\_manifest\_\_.py})}

\begin{center}
\small
\begin{tabularx}{\textwidth}{lX}
\toprule
\textbf{Nhóm} & \textbf{Module} \\
\midrule
Nền Wujia & \texttt{wujia\_core} \\
Nghiệp vụ sạch & \texttt{wujia\_franchise}, \texttt{wujia\_franchise\_contract},
  \texttt{wujia\_franchise\_inspection}, \texttt{wujia\_franchise\_operations},
  \texttt{wujia\_sale}, \texttt{wujia\_fleet}, \texttt{wujia\_delivery}, \texttt{wujia\_account} \\
Khung portal & \texttt{wujia\_portal\_layout} (chỉ depend \texttt{base, web, auth\_signup}),
  \texttt{wujia\_portal\_base} (depend \texttt{wujia\_sale} + layout) \\
Portal thuần (controller + QWeb) & \texttt{portal\_debt}, \texttt{portal\_delivery},
  \texttt{portal\_inspection}, \texttt{portal\_purchase\_history}, \texttt{portal\_report},
  \texttt{portal\_sale} \\
\textbf{``Portal'' nhưng ôm nghiệp vụ} & \texttt{portal\_order\_window}, \texttt{portal\_info\_request},
  \texttt{portal\_knowledge}, \texttt{portal\_support}, \texttt{portal\_notification},
  \texttt{portal\_exam}, \texttt{portal\_return} \\
Workstream riêng & \texttt{wj\_ks\_dashboard\_ninja}, \texttt{wj\_ks\_dn\_advance} \\
\bottomrule
\end{tabularx}
\end{center}

Bên thứ ba nằm cùng \texttt{custom/} (kế toán \texttt{om\_*}/\texttt{eh\_*},
\texttt{auditlog}, \texttt{base\_user\_role}, \texttt{purchase\_request}, \texttt{mcp\_server},
\texttt{legion\_enterprise\_theme}\ldots) và theme backend MuK trong \texttt{themes/} được
coi là tầng Odoo.

\subsection*{Chẩn đoán}

\paragraph{(a) Chỗ chủ dự án lo lại đang sạch.} Không module nghiệp vụ nào
(\texttt{core}, \texttt{franchise*}, \texttt{sale}, \texttt{fleet}, \texttt{delivery},
\texttt{account}) depend bất kỳ \texttt{wujia\_portal\_*} nào. Do đó
\texttt{wujia\_mobile\_franchise} $=$ mobile core $+$ \texttt{wujia\_franchise}
\textbf{không kéo theo portal}.

\paragraph{(b) Chỗ hỏng thật: 7 module ``portal'' đang ôm nghiệp vụ backend.} Chúng định nghĩa
model, view backend, menu, nhóm quyền, sequence. Tức là \emph{nghiệp vụ nằm trong kênh}.

\begin{center}
\footnotesize
\begin{tabular}{lrrrrl}
\toprule
Module & xmlid & Dòng model & Dòng ACL & Sequence & Model sở hữu \\
\midrule
\texttt{portal\_order\_window}  & 7  & 291  & 3  & 0 & order.window (+SO) \\
\texttt{portal\_info\_request}  & 13 & 175  & 3  & 1 & info.update.request \\
\texttt{portal\_knowledge}      & 20 & 250  & 7  & 1 & article / category / tag \\
\texttt{portal\_support}        & 26 & 310  & 5  & 1 & ticket / category \\
\texttt{portal\_notification}   & 29 & 566  & 10 & 1 & notification / read / type \\
\texttt{portal\_exam}           & 35 & 646  & 13 & 1 & 5 model thi \\
\texttt{portal\_return}         & 32 & 1108 & 12 & 1 & return, compensation (+SO/picking) \\
\bottomrule
\end{tabular}
\end{center}

Hệ quả cụ thể: muốn cho HQ duyệt \emph{Đăng ký thi} trên điện thoại thì
\texttt{wujia\_mobile\_exam} buộc phải depend \texttt{wujia\_portal\_exam}, kéo theo
\texttt{wujia\_portal\_base}, \texttt{wujia\_portal\_layout}, \texttt{wujia\_sale} và cả bundle
\texttt{web.assets\_frontend}. Khi tách repo cũng không có cách nào để nghiệp vụ thi nằm ngoài
folder portal (Hình~\ref{fig:adr27-why}).

\adrWhyFigure

\paragraph{(c) Ba điểm lệch nhỏ hơn, ghi lại cho đủ.}
\begin{enumerate}
\item \texttt{wujia\_portal\_base} đọc \emph{mềm} model của module nó không depend
  (\texttt{request.env.get('wujia.knowledge.article')}, \texttt{'wujia.notification.read'}\ldots)
  để dựng Home. Không lỗi khi thiếu module, nhưng lõi portal đang ``biết'' tên model của nghiệp vụ.
\item Test component của \texttt{wujia\_portal\_layout} (\texttt{test\_d3/d4/e2b/c8\ldots})
  tham chiếu template của \texttt{portal\_exam}, \texttt{portal\_return}\ldots\ Tầng khung đang
  test ngược lên tầng trên.
\item \texttt{wujia\_franchise\_inspection} chứa một trang web cho người đi khảo sát
  (\texttt{/franchise/inspection/do}, \texttt{auth='user'}) cùng CSS/JS backend. Đây là một màn
  \emph{kênh} nằm trong module nghiệp vụ, và cũng là ứng viên tự nhiên nhất cho mobile.
\end{enumerate}

\paragraph{(d) Hai bundle CSS vốn đã tách sẵn.} Portal nạp vào \texttt{web.assets\_frontend}, backend
nạp vào \texttt{web.assets\_backend}. Theme backend đang chạy trên UAT là
\texttt{muk\_web\_theme} 19.0.1.4.4 (\texttt{legion\_enterprise\_theme} có trong repo nhưng
\emph{uninstalled}). Nên ``mobile core chỉ sửa CSS backend, không ảnh hưởng portal'' là khả thi
về kỹ thuật.

\paragraph{(e) Đã có tiền lệ tách module thành công.} Đợt tách \texttt{wujia\_franchise}
$\rightarrow$ \texttt{\_contract}/\texttt{\_inspection}/\texttt{\_operations} dùng
\texttt{pre\_init\_hook} đổi chủ \texttt{ir\_model\_data}
(\texttt{wujia\_franchise\_inspection/hooks.py}), và \textbf{đã cài trên UAT}
(\texttt{wujia\_franchise\_inspection} 19.0.1.0.0 \emph{installed}). Quy trình đề xuất ở dưới
chính là quy trình đó, chuẩn hoá lại.

\section{Quyết định --- 4 tầng, mũi tên chỉ đi xuống}

\textbf{Nói ngắn:} nghiệp vụ (\texttt{wujia\_exam}\ldots) và khung giao diện
(\texttt{portal\_layout}, \texttt{mobile\_core}) là \textbf{hai thứ ngang hàng, không biết nhau}.
Chỉ module ghép ở tầng trên (\texttt{portal\_exam}, \texttt{mobile\_exam}) mới depend cả hai.
Sơ đồ ở Hình~\ref{fig:adr27-layers}.

\begin{itemize}
\item \textbf{L0} Odoo + bên thứ 3 (\texttt{base, web, sale, stock, account, portal, muk\_web\_theme, om\_*}\ldots).
\item \textbf{L1} nền: \texttt{wujia\_core}; platform \texttt{wujia\_security/approval/job/audit/\ldots\_core} và \texttt{wujia\_ui\_core} --- \textbf{chỉ tạo khi thật sự có ít nhất 2 nơi cần} (xem \S\ref{sec:adr27-ba}).
\item \textbf{L2} gồm hai nhóm \textbf{cùng cấp}:
  \begin{itemize}
  \item nghiệp vụ: 8 module đang sạch + 7 phân hệ tách ra (\texttt{wujia\_order\_window}, \texttt{\_info\_request}, \texttt{\_knowledge}, \texttt{\_support}, \texttt{\_announcement}, \texttt{\_exam}, \texttt{\_return});
  \item khung kênh: \texttt{wujia\_portal\_layout} \quad|\quad \texttt{wujia\_mobile\_core}.
  \end{itemize}
\item \textbf{L3} module ghép, chia hai bậc:
  \begin{itemize}
  \item \textbf{L3a} \texttt{wujia\_portal\_base} --- nền ghép của nhánh portal (``tầng 2.5'');
  \item \textbf{L3b} \texttt{wujia\_portal\_<x>} \quad|\quad \texttt{wujia\_mobile\_<x>}.
  \end{itemize}
\end{itemize}

\adrLayerFigure

\subsection{Bốn luật dependency}

\begin{enumerate}
\item \textbf{Nghiệp vụ không depend khung kênh hay module ghép.} Nghiệp vụ không biết portal hay
  mobile tồn tại. Model, view backend, menu, nhóm quyền, ACL, ir.rule (kể cả rule cho
  \texttt{group\_portal}), sequence, cron \textbf{đều ở nghiệp vụ}.
\item \textbf{Khung kênh không depend nghiệp vụ.} \texttt{wujia\_portal\_layout} chỉ depend
  \texttt{base, web, auth\_signup} (đúng như code hiện tại); \texttt{wujia\_mobile\_core} chỉ depend
  \texttt{web, muk\_web\_theme}. Cả hai phải bê sang dự án khác được.
\item \textbf{Hai kênh không depend chéo.} \texttt{wujia\_mobile\_*} không depend
  \texttt{wujia\_portal\_*} và ngược lại. Thứ gì cả hai cần chung thì mới tạo \texttt{wujia\_ui\_core}.
\item \textbf{Module ghép chỉ chứa phần kênh.} Portal: controller, QWeb, frontend asset. Mobile:
  backend asset, view inherit tối thiểu. \textbf{Không định nghĩa model nghiệp vụ ở đây.} Model thuần
  phục vụ kênh được ở lại, ví dụ giỏ hàng \texttt{wujia.portal.cart}, seam \texttt{wujia.portal.debt}.
\end{enumerate}

\subsection{\texttt{wujia\_portal\_base} --- nền của nhánh portal (ngoại lệ có ghi nhận)}

\texttt{wujia\_portal\_base} là \texttt{portal\_layout} + cửa hàng + Home, nên nó đứng
ở \textbf{L3a --- bậc dưới của tầng ghép}, không phải khung kênh. Mọi \texttt{wujia\_portal\_<x>} depend nó.

\begin{itemize}
\item \textbf{Giữ depend \texttt{wujia\_franchise} + \texttt{wujia\_sale}.} Portal Wujia là kênh đặt
  hàng của cửa hàng: không có cảnh chạy portal mà thiếu cửa hàng hay thiếu bán hàng.
\item \textbf{Luật L3a/L3b} (kiểm bằng \texttt{check\_layers.py}): mọi \texttt{wujia\_portal\_<x>}
  phải depend \texttt{wujia\_portal\_base}; \texttt{wujia\_portal\_base} không được depend bất kỳ
  \texttt{wujia\_portal\_<x>} nào.
\item \textbf{Cấm thêm depend nghiệp vụ mới.} Phân hệ khác chỉ được đọc mềm (\texttt{request.env.get(\ldots)}) như hiện tại.
\item \textbf{Không tách Home ra từng module lúc này.} Home vừa gom về một chỗ (C7, E1). Rải khối Home
  ra nhiều module qua xpath làm thứ tự khối phụ thuộc thứ tự cài, dễ vỡ mà không được gì.
\item \textbf{Mobile không cần module tương đương.} Backend đã có record rule lo quyền dữ liệu, nên
  \texttt{mobile\_core} giữ được thuần. Nếu sau này mobile cần chọn cửa hàng thì làm
  \texttt{wujia\_mobile\_franchise} ở tầng ghép.
\end{itemize}

\subsection{Vì sao chọn cấu trúc này}

\begin{itemize}
\item \textbf{Đúng cách Odoo tự tổ chức}: \texttt{sale} (nghiệp vụ), \texttt{website\_sale}
  (kênh web), \texttt{sale\_stock} (cầu nối). Ai quen Odoo đọc là hiểu, không cần học quy ước riêng.
\item \textbf{Tách repo chỉ là cắt theo tầng}: \texttt{wujia-core/} (L1+L2),
  \texttt{wujia-portal/} (khung + ghép portal), \texttt{wujia-mobile/} (khung + ghép mobile). Mỗi folder là
  một mục \texttt{addons\_path}. Hai khung chung chung (\texttt{portal\_layout},
  \texttt{mobile\_core}) còn có thể lên kho dùng chung beesuite.
\item \textbf{Cài theo nhu cầu}: môi trường không cần portal chỉ cài L1 + nghiệp vụ (+ mobile), backend
  vẫn đủ.
\item \textbf{Luật đọc được bằng máy}: chỉ cần đọc \texttt{\_\_manifest\_\_.py} là kiểm được. Sẽ
  có script \texttt{scripts/qa/check\_layers.py} báo vi phạm, để luật không chỉ nằm trên giấy.
\end{itemize}

\subsection{Phương án đã cân nhắc và loại}

\begin{description}
\item[Giữ nguyên, chỉ áp luật cho module mới.] Loại. Mobile cho thi, thông báo, hỗ trợ, đổi trả
  gần như chắc chắn sẽ cần (HQ duyệt trên điện thoại). Đến lúc đó hoặc phá luật 3, hoặc tách giữa
  lúc đang làm tính năng, tức là trả giá gấp đôi.
\item[Tách cả 7 phân hệ trong một đợt.] Loại. 7 lần đổi chủ xmlid trên UAT có dữ liệu thật trong
  một lệnh deploy thì hỏng không bisect được. Sai một xmlid là mất nhóm quyền hoặc sequence một
  cách \emph{im lặng}. Bài học D3c/D4: lỗi deploy chỉ lộ khi đo.
\item[Mọi phân hệ luôn đủ bộ 3 module.] Loại. Module rỗng làm phình danh sách mà không đem lại gì.
\end{description}

\section{Mobile core --- phụ thuộc và phạm vi}

\textbf{Quyết định:} \texttt{wujia\_mobile\_core} depend \texttt{web} $+$ \texttt{muk\_web\_theme},
\textbf{không} depend \texttt{wujia\_core}.

\begin{itemize}
\item \texttt{wujia\_core} chỉ có bảng khu vực/phường. CSS mobile không dùng tới. Depend vào đó
  thì mobile core thành ``của Wujia'', không tái dùng được, và phá luật 3.
\item \texttt{muk\_web\_theme} là nền backend \emph{đang thật sự render} (appsbar, chatter,
  dialog). CSS mobile viết đè lên DOM của MuK, nên phải nạp \emph{sau} MuK. Không depend thì thứ
  tự nạp asset không được bảo đảm, và rule sẽ thua đúng kiểu bẫy C6/D3 (rule nạp sau thắng).
\item Nếu sau này đổi theme backend (ví dụ sang \texttt{legion\_enterprise\_theme}), chỉ một mình
  mobile core phải sửa. Không phân hệ nào bị kéo theo.
\end{itemize}

\textbf{Nội dung:} chỉ bundle \texttt{web.assets\_backend}, bọc trong breakpoint mobile.
\begin{itemize}
\item Trạng thái chạm \texttt{:active} thay cho \texttt{:hover}; vùng chạm tối thiểu 44px.
\item Chọn bản ghi trên kanban/list theo kiểu app (chạm giữ / ô chọn hiện rõ).
\item Áp cho các view \emph{chung}: kanban, list, form, control panel, search panel, dialog,
  chatter. Dự kiến gánh phần lớn nhu cầu.
\item \textbf{Không} model Python, \textbf{không} controller, \textbf{không} chạm
  \texttt{web.assets\_frontend}. Portal không thể bị ảnh hưởng.
\end{itemize}

\section{Module cầu nối \texttt{wujia\_mobile\_<x>} --- \texttt{auto\_install}, chỉ tạo khi có nội dung}

\textbf{Cầu nối là gì.} Module nhỏ chỉ để nối hai module đã có. Odoo có \texttt{sale\_stock}:
depend \texttt{['sale', 'stock']} và khai \texttt{auto\_install: True}. Nghĩa là khi DB
\textbf{đã có đủ} \texttt{sale} lẫn \texttt{stock} thì Odoo \textbf{tự cài}
\texttt{sale\_stock}, không ai phải nhớ. Áp vào Wujia:

\begin{lstlisting}[style=python]
# wujia_mobile_franchise_inspection/__manifest__.py
'depends': ['wujia_mobile_core', 'wujia_franchise_inspection'],
'auto_install': True,
'assets': {'web.assets_backend': [...]},
\end{lstlisting}

\textbf{Khi nào tạo.} Chỉ khi CSS/view \textbf{nhắc tới model hoặc view cụ thể của một phân hệ}.
Ví dụ phiếu khảo sát cửa hàng, hay form duyệt đổi trả nhiều cột. Còn view chung đã đẹp nhờ mobile
core thì \emph{không tạo}. Một selector kiểu \texttt{.o\_wujia\_inspection\_*} nằm trong mobile
core là dấu hiệu phải tách ra cầu nối.

Tính đối xứng nằm ở \textbf{luật đặt tên} (\texttt{wujia\_<x>} / \texttt{wujia\_portal\_<x>} /
\texttt{wujia\_mobile\_<x>}), không nằm ở việc lúc nào cũng phải có đủ ba module.

\section{Phác thảo ERD theo module đích}

Hình~\ref{fig:adr27-erd} xếp model chính vào \emph{module mà nó sẽ thuộc về} sau khi tách, cùng
quan hệ Many2one then chốt. Danh sách đo bằng AST trên \texttt{custom/*/models} (không đoán theo
tên). Ba điều đọc ra từ hình:

\begin{enumerate}
\item \textbf{Mọi đường đều dẫn về \texttt{wujia.franchise.management}.} Đây là hub. Nên
  \texttt{wujia\_franchise} phải đứng thấp nhất trong L2, và mọi phân hệ khác depend nó chứ không
  có chiều ngược lại.
\item \textbf{Quan hệ ngược thì dùng kế thừa từ module trên.} \texttt{wujia\_franchise\_contract}
  thêm \texttt{current\_contract\_id} vào management, \texttt{\_inspection} thêm
  \texttt{latest\_inspection\_id}. Hub không phải biết hai module đó. Đây là cách tránh vòng lặp
  phụ thuộc mà ADR-015 từng vướng, và là khuôn cho mọi phân hệ tách ra.
\item \textbf{Depends của phân hệ tách ra suy thẳng từ quan hệ}: \texttt{wujia\_return},
  \texttt{wujia\_support} $\rightarrow$ \texttt{wujia\_sale} (\texttt{sale.order},
  \texttt{stock.picking.batch}); \texttt{wujia\_order\_window} $\rightarrow$ \texttt{wujia\_sale}
  $+$ \texttt{wujia\_core}; \texttt{wujia\_exam}, \texttt{wujia\_announcement},
  \texttt{wujia\_info\_request} $\rightarrow$ \texttt{wujia\_franchise};
  \texttt{wujia\_knowledge} không có Many2one nào sang Wujia, chỉ cần \texttt{mail} (hiện depend
  \texttt{wujia\_franchise} có lẽ vì quyền, sẽ kiểm lúc tách).
\end{enumerate}

\adrErdFigure

\section{Quy trình tách một phân hệ ra khỏi module portal}
\label{sec:adr27-xmlid}

\textbf{Điều gì làm nó nguy hiểm.} Bảng DB \emph{không} đổi, vì tên bảng lấy từ \texttt{\_name}
chứ không lấy từ tên module. Cái nguy hiểm là \texttt{ir\_model\_data}: mọi view, menu, nhóm quyền,
rule, ACL, sequence, \texttt{ir.model}, \texttt{ir.model.fields} đang ghi
\texttt{module='wujia\_portal\_<x>'}. Nếu chỉ chuyển file, lúc \texttt{-u} Odoo sẽ tạo \emph{bản
sao} dưới tên module mới, rồi \texttt{\_process\_end} \textbf{dọn bản cũ}. Hệ quả: user rơi khỏi
nhóm quyền, rule biến mất, sequence mã phiếu chạy lại từ đầu, tệ nhất là rơi dòng
\texttt{ir.model}.

\textbf{Các bước} (đã chạy thật ở pilot F7 --- \texttt{wujia\_portal\_order\_window} $\rightarrow$
\texttt{wujia\_order\_window}, 25/09/2026; F8--F13 làm theo đúng khuôn này):
\begin{enumerate}
\item \textbf{Chụp mốc trước} trên DB copy cô lập bằng \texttt{scripts/qa/split\_snapshot.py}
  (\texttt{--modules} cũ,mới \texttt{--models} \ldots\ \texttt{--icp} \ldots\ \texttt{--seq} \ldots).
  DB copy phải có dữ liệu thật của phân hệ; không có thì seed trước, vì bảng rỗng thì phép đo
  Pass rỗng (bài học D6d). Nhãn dịch chỉ đo được khi DB đã bật \texttt{vi\_VN}.
\item Tạo \texttt{wujia\_<x>} bằng \texttt{git mv}: \texttt{models/}, view backend, menu,
  \texttt{security/}, \texttt{data/} (sequence, cron), \texttt{i18n/vi\_VN.po}. Trong
  \texttt{.po} phải đổi cả tiền tố xmlid (\texttt{model:ir.model.fields,\ldots:wujia\_portal\_<x>.field\_\ldots}
  $\rightarrow$ \texttt{wujia\_<x>.}), không thì bản dịch không gắn được vào xmlid mới và nhãn
  rơi về tiếng Anh mà không có dòng log nào.
  \textbf{Bẫy data \texttt{noupdate} (F9):} lần \texttt{-i} module mới chạy chế độ \texttt{init}, Odoo
  \emph{nạp lại} mọi bản ghi \texttt{noupdate="1"} lên xmlid vừa đổi chủ (\texttt{convert.py}: chỉ bỏ qua
  noupdate khi \texttt{mode != 'init'}). Dòng \texttt{<field name="number\_next">1</field>} trong XML
  sequence $\Rightarrow$ \texttt{ir.sequence.write} $\Rightarrow$ \texttt{ALTER SEQUENCE RESTART 1}:
  mã chứng từ chạy lại từ 1 và trùng mã cũ. \textbf{Xoá dòng \texttt{number\_next} khỏi XML} trước khi
  tách (mặc định vẫn là 1 khi cài mới). Cron, tham số hệ thống, mail template noupdate cũng bị ghi đè
  bằng giá trị XML $\Rightarrow$ đo giá trị trên UAT trước deploy, lệch với XML thì dừng hỏi.
\item \textbf{Model dính chặt nhau đi cùng một module} (One2many $\leftrightarrow$ Many2one hai
  chiều). Chiều ngược sang module khác thì thêm bằng \texttt{\_inherit} ở module phía trên (bài
  học ADR-015).
\item \texttt{hooks.py} + \texttt{'pre\_init\_hook'} \textbf{khai trong manifest và import ở
  \texttt{\_\_init\_\_.py}}. Hook khuôn Khảo sát hiện không còn được nối (gỡ khỏi manifest ở
  \texttt{ec6d380}, 16/09), nên chép file thôi thì hook không chạy. Tiện ích dùng chung nằm ở L1:
  \texttt{from odoo.addons.wujia\_core.tools.module\_split import imd\_names, migrate\_ownership}
  (FR-P dời từ \texttt{wujia\_order\_window/hooks.py}). \texttt{migrate\_ownership(cr, old, new,
  names, models)} đổi chủ ba chỗ (và từ chối chạy khi \texttt{new} chưa có trong
  \texttt{ir\_module\_module} --- không thì ràng buộc bị gán chủ NULL trong im lặng):
  \begin{itemize}
  \item \texttt{ir\_model\_data}: \texttt{names=None} chuyển \emph{tất cả} (module cũ bị rút hết,
    như F7); tách một phần (F8+) thì \texttt{names=imd\_names(cr, old, models, extra=[menu, group,
    dữ liệu mẫu])} --- helper suy xmlid của model/field/selection/inherit/constraint/access/rule/
    view/action/server action/cron/mail.template/sequence gắn với \texttt{models} (kể cả model kế
    thừa như \texttt{sale.order}, chỉ dòng do \texttt{old} sở hữu). Menu/group không có cột model
    nên liệt kê tay qua \texttt{extra}. Luôn in phần dư \texttt{SELECT name FROM ir\_model\_data
    WHERE module=old} trừ kết quả helper: FR-P thử trên \texttt{info\_request} ra 72/75, ba dòng dư
    đúng là ba QWeb template portal phải ở lại; đối chứng F7 36/36.
  \item \texttt{ir\_model\_constraint} và \texttt{ir\_model\_relation}: cột \texttt{module} là
    khoá ngoài tới \texttt{ir\_module\_module}, \emph{không} nằm trong \texttt{ir\_model\_data}.
    Khuôn Khảo sát bỏ sót chỗ này. Đo mutation ở F7: bỏ bước này thì \textbf{không} mất ràng buộc
    (Odoo gỡ theo xmlid, và \texttt{unlink} tự bỏ qua khi còn dòng cùng tên), nhưng module mới
    sinh 2 dòng CHECK trùng tên và 6 dòng cũ mồ côi sau khi gỡ vỏ.
  \end{itemize}
  Chạy ở \texttt{pre\_init\_hook} vì Odoo \emph{không} chạy migration script ở lần cài đầu.
\item Module cũ: còn controller/QWeb/asset thì giữ, thay depends nghiệp vụ bằng
  \texttt{wujia\_<x>}, xoá \texttt{models/}, \texttt{security/}, \texttt{data/} đã rỗng. \texttt{.po}
  tách theo dòng tham chiếu \texttt{\#:} (controller + QWeb portal ở lại; một msgid dùng ở cả hai
  phía thì ghi ở cả hai file). F8: 137 mục thành 56 portal + 81 nghiệp vụ; F9: 123 thành 26 + 97. Rút hết
  (như F7) thì để \textbf{vỏ rỗng} (\texttt{\_\_init\_\_.py} trống, \texttt{data: []}, depends
  mỗi \texttt{wujia\_<x>}) và thêm vào \texttt{DEPRECATED} của \texttt{check\_layers}. Bump version
  cả hai.
\item Sửa tham chiếu \texttt{wujia\_portal\_<x>.\ldots} ở module khác, cùng manifest và
  \texttt{import odoo.addons.wujia\_portal\_<x>} (F7: \texttt{portal\_sale} đổi depends +
  import lớp lỗi \texttt{OrderWindowClosed}), \texttt{reseed\_full.sh/.ps1},
  \texttt{.github/workflows/deploy.yml} (danh sách \texttt{-i}/\texttt{-u}; FR-P phát hiện F7 bỏ
  sót). \textbf{Grep cả module mobile của anh Thái}: \texttt{wujia\_mobile\_portal\_info\_request}
  và \texttt{\_exam} ghi đè action/\texttt{ref()} view bằng xmlid \texttt{wujia\_portal\_<x>.*} ---
  đổi chủ xmlid xong thì \texttt{ref} ở đó gãy lúc \texttt{-u}. \textbf{Chốt F8 (26/09): sửa nguồn
  mobile} (đổi tiền tố + depends), không giữ xmlid cũ. Đo F8 trên DB copy có cài module mobile, khi
  chưa sửa: \texttt{-u} dừng ở \texttt{Cannot update missing record
  'wujia\_portal\_info\_request.action\_\ldots'}, exit 255. \textbf{Không phải huỷ cả lượt}: Odoo 19
  commit sau từng module, nên module mới đã cài và đổi chủ xong, module portal đã lên version mới;
  chỉ module mobile dừng lại. Khởi động thường vẫn chạy, action mobile vẫn đủ list/kanban/form.
  Nhưng mọi lần \texttt{-u} sau vẫn gãy cho tới khi sửa nguồn mobile.
\item \textbf{Deploy một lệnh:} \texttt{-i wujia\_<x> -u wujia\_portal\_<x>,<module tham chiếu>}.
  Đo F7: chỉ \texttt{-u wujia\_portal\_<x>} (quên \texttt{-i}) cũng tự cài module mới, và
  \texttt{update\_list} cập nhật luôn bảng phụ thuộc từ manifest mới.
\item \textbf{Chụp mốc sau} rồi so: \texttt{split\_snapshot.py --diff trước sau --rename
  wujia\_portal\_<x>=wujia\_<x>}. Kỳ vọng \emph{chỉ cột module đổi}; mọi dòng ``LỆCH THẬT''
  phải giải trình được (F7: 1 dòng --- view Settings đổi \texttt{name} app/block theo tên module
  mới). Từ FR-P snapshot đo thêm \texttt{inherit} (\texttt{ir\_model\_inherit}, Odoo 19 có xmlid
  riêng), \texttt{sel} (giá trị selection), \texttt{cron}/\texttt{tpl}/\texttt{srv},
  \texttt{attach}, rule theo module --- F7 không gặp vì model không kế thừa mixin, F8 gặp ngay
  (\texttt{info\_request}: 2 inherit + 14 selection). Route controller không nằm trong DB: kiểm
  bằng \texttt{controller\_inventory.py}.
\item \textbf{Gỡ vỏ} (khi module cũ rút hết): đo trên DB copy trước --- \texttt{button\_immediate\_uninstall}
  rồi chụp lần ba, phải 0 lệch. Trên UAT: sau deploy, đo module cũ còn \textbf{0 xmlid} rồi
  Uninstall một cú trên Apps (FR-P làm qua RPC \texttt{button\_immediate\_uninstall}, 6s, chụp
  lại 0 lệch); xoá thư mục khỏi repo ở phiên review khi mọi DB đã gỡ. Dòng
  \texttt{ir\_module\_module} của vỏ vẫn còn (\texttt{uninstalled}) trên DB đã cài, nhưng DB trắng
  thì không --- test nào tra \texttt{id} của vỏ sẽ ngã (FR-P sửa \texttt{test\_split\_ownership}
  dùng dòng module giả).
\end{enumerate}

\textbf{Kiểm bằng số, trước và sau, trên DB copy cô lập} (không đụng \texttt{wujia\_tea\_19}):
\texttt{split\_snapshot} gom trong một lượt: \texttt{ir\_model\_data} theo module; số dòng + md5
toàn bảng; cột và ràng buộc Postgres thật; \texttt{ir\_model\_constraint}/\texttt{relation} gom
theo \emph{tên} (để dòng trùng tên hiện ra, không gộp mất); field (kiểu, nhãn en/vi); ACL; rule;
đường dẫn menu; action; view (md5 arch); \texttt{ir\_config\_parameter}; số user mỗi nhóm quyền;
\texttt{number\_next} của sequence và \texttt{last\_value} của sequence Postgres (kiểu
\texttt{standard} chỉ có số thật ở đó --- cột \texttt{number\_next} luôn là 1, F9 mới bổ sung). Kèm: test cũ xanh y hệt; DB trắng chỉ cài \texttt{wujia\_<x>}
chạy được test của nó; HTML các màn liên quan HEAD $\leftrightarrow$ mới giống từng byte (chuẩn
hoá hash asset + \texttt{registry\_hash}).

\textbf{Số đo pilot F7:} 42 dòng đổi chủ (36 xmlid + 6 ràng buộc), 1 lệch có giải trình, gỡ vỏ
0 lệch, 0 ERROR khi nâng cấp; suite 833 test 0 đỏ; DB trắng 10/10; HTML 5/5 route giống từng
byte; mutation 5/5 đỏ. \textbf{FR-P (26/09):} UAT gỡ vỏ xong, đo lại 0 lệch; thư mục vỏ xoá;
\texttt{check\_layers} 0 vi phạm phía Dev; thử tách một phần \texttt{info\_request} trên DB copy:
80 dòng đổi chủ, 0 lệch thật, sequence giữ nguyên.

\textbf{F8 (26/09) --- tách một phần \texttt{info\_request} làm thật:} hook
\texttt{imd\_names} + \texttt{extra=[menu]} chuyển 72 xmlid + 7 ràng buộc + 1 quan hệ; module portal
còn đúng 3 QWeb; snapshot 80 dòng đổi chủ, \textbf{0 lệch thật}, 0 ERROR. Controller mỏng: gate
Owner/Manager, phạm vi cửa hàng và ``tạo + đính kèm + gửi'' về model
(\texttt{\_portal\_can\_request}, \texttt{\_portal\_scope\_domain}, \texttt{create\_from\_portal}
trong một savepoint; tiện ích upload vẫn ở \texttt{portal\_base}, truyền vào qua callback nên L2
không import portal). HTML 5 route + JSON giống từng byte; mutation 3/3 đỏ. \textbf{Bẫy test:}
\texttt{assertRaises} của Odoo tự bọc savepoint, nên test ``lỗi thì rollback nguyên khối'' viết bằng
nó sẽ Pass cả khi code quên savepoint. Phải dùng \texttt{try/except} thường (bắt được qua mutation).

\textbf{F9 (26/09) --- tách một phần \texttt{knowledge}:} 3 model, 4 menu, cron, sequence KNW- sang
\texttt{wujia\_knowledge}; hook đổi chủ 107 xmlid + 13 ràng buộc + 2 quan hệ, portal còn 3 QWeb + 2 mục
nav. Lượt deploy đầu trên DB copy \textbf{reset sequence KNW- từ 438 về 1} mà snapshot cũ báo 0 lệch
(chỉ đọc \texttt{number\_next}) --- bẫy \texttt{noupdate} ở bước 2; sửa XML + snapshot, chạy lại
\textbf{0 lệch thật}. UAT lúc đó có 27 bài, số kế 28: không sửa thì bài kế ra \texttt{KNW-000001}
(trùng). F8 cùng bẫy nhưng vô hại: UAT 0 yêu cầu, \texttt{INF-} số kế 1. Controller mỏng: luật
``bài nào hiện trên portal'', tìm kiếm, ``đính kèm thuộc bài'' về model
(\texttt{\_portal\_visible\_domain}, \texttt{\_portal\_search\_domain}, \texttt{\_portal\_get\_attachment});
Home (\texttt{portal\_base}, không depend) gọi chung luật qua \texttt{hasattr} --- vá lệch cũ: Home
không lọc ngày phát hành nên bài hẹn giờ đã hiện ở Home. HTML 12/12 route + JSON giống từng byte, Home
lệch đúng bài hẹn giờ. Mutation: 2 luật model đỏ 2/2 test mỗi luật; hook quên \texttt{extra} thì test
đơn vị \emph{không} bắt (trạng thái cuối giống hệt, chỉ khác id menu), snapshot bắt 10 lệch
--- lỗi hook là việc của phép đo trước/sau, không thay bằng test được.

\section{Lộ trình}

\begin{center}
\small
\begin{tabularx}{\textwidth}{rlX}
\toprule
\# & Phiên & Lý do thứ tự \\
\midrule
P1 & Dựng \texttt{wujia\_mobile\_core} & Không phụ thuộc việc tách, chạy song song được. Cần BA chốt màn ưu tiên. \\
P2 & Tách \texttt{order\_window} (pilot) & \textbf{Xong F7 (25/09)}: 36 xmlid + 6 ràng buộc đổi chủ, 0 nhóm quyền; module cũ thành vỏ chờ gỡ. Lộ ra bẫy \texttt{ir\_model\_constraint} và hook khuôn không còn được nối. \\
P3 & \texttt{info\_request} & \textbf{Xong F8 (26/09)}: tách một phần, 72 xmlid + 7 ràng buộc + 1 quan hệ đổi chủ, 0 nhóm quyền; 3 QWeb + controller mỏng ở lại portal. \\
P4 & \texttt{knowledge} & \textbf{Xong F9 (26/09)}: tách một phần, 107 xmlid + 13 ràng buộc + 2 quan hệ đổi chủ, cron + sequence giữ nguyên; lộ bẫy data \texttt{noupdate} reset sequence. \\
P5 & \texttt{support} & 26 xmlid, nhiều depends (\texttt{sale}, \texttt{stock\_picking\_batch}, \texttt{sales\_team}). \\
P6 & \texttt{announcement} & Tách từ \texttt{portal\_notification}. 2 nhóm quyền, sequence ANN, ir.rule theo cửa hàng. Giữ \texttt{\_name} \texttt{wujia.notification} (không đổi tên bảng). \\
P7 & \texttt{exam} & 2 nhóm quyền, 5 model. \\
P8 & \texttt{return} & Lớn nhất: 1108 dòng model, kế thừa SO/picking, hook picking, wizard SO 0đ FIFO. Để cuối, khi quy trình đã chín. \\
\bottomrule
\end{tabularx}
\end{center}

Nếu BA cần mobile cho một phân hệ sớm thì \textbf{đẩy phân hệ đó lên}. Riêng pilot
\texttt{order\_window} vẫn đi trước vì nó rẻ. Ba điểm lệch nhỏ ở mục (c) được dọn kèm phiên
tách phân hệ tương ứng, không mở phiên riêng.

\section{Đối chiếu với bảng BA ``0. Kiến trúc module''}
\label{sec:adr27-ba}

BA có một bảng kiến trúc ở đầu tab \emph{1. Model/ Field} (dòng 1--21). Đã đọc từng dòng và đối
chiếu với manifest trong code (\texttt{12750a2}) và trạng thái cài trên UAT.

\subsection*{BA đang chia thế nào}

Hai cấp. \textbf{Tầng 1} gom ba loại khác nhau: platform core dùng chung (\texttt{ui, mobile,
security, approval, notification, document, integration, job, audit}), nghiệp vụ
(\texttt{franchise, sale, fleet, delivery}) và lõi kênh/báo cáo (\texttt{portal\_base, report}).
\textbf{Tầng 2} là module con ``\,$\hookrightarrow$ thuộc'' một module Tầng 1.

\subsection*{Kết luận: khung theo ADR-027, lấy thêm 4 ý từ BA}

Khung tầng giữ nguyên, vì bảng BA không tách platform / nghiệp vụ / kênh nên không đặt được
luật ai được depend ai, tức không giữ được hai kênh tách nhau. Bốn ý \textbf{lấy vào}:

\begin{enumerate}
\item \textbf{Tầng platform core nằm ngang} dưới nghiệp vụ (đã thêm vào L1). Điều kiện: chỉ tạo khi
  có ít nhất 2 phân hệ cần, và \textbf{bọc module đã cài trên UAT} thay vì viết lại:
  \texttt{queue\_job} (job), \texttt{auditlog} (audit), \texttt{dynamic\_approval\_workflow} +
  \texttt{eh\_account\_approval} (duyệt).
\item \textbf{\texttt{wujia\_ui\_core}}: chỗ để portal và mobile dùng chung giao diện mà không depend
  nhau (đã thêm vào luật 2).
\item \textbf{Ký hiệu ``$\hookrightarrow$ thuộc module X''} cho module con, dễ đọc với BA. Ví dụ
  \texttt{franchise\_contract/\_inspection/\_operations} $\hookrightarrow$ \texttt{franchise};
  \texttt{portal\_<x>} $\hookrightarrow$ \texttt{portal\_base}; \texttt{mobile\_<x>}
  $\hookrightarrow$ \texttt{mobile\_core}.
\item \textbf{Đổi tên phân hệ thông báo khi tách thành \texttt{wujia\_announcement}}, vì BA đã dành
  tên \texttt{notification\_core} cho nền tảng gửi thông báo (sự kiện + kênh). Chỉ đổi tên
  \emph{module}; \texttt{\_name} vẫn là \texttt{wujia.notification} để không phải đổi tên bảng.
\end{enumerate}

\subsection*{Chỗ ADR-027 khác BA}

\begin{itemize}
\item BA xếp \texttt{portal\_return} là ``Tầng 2 --- thuộc Portal'' với mô tả gồm cả \emph{duyệt/xử
  lý}. ADR-027 tách phần duyệt/xử lý xuống \texttt{wujia\_return} (L2). Đây là \textbf{câu hỏi số 6}
  bên dưới.
\item BA ghi \texttt{mobile\_core} depend \texttt{wujia\_core + web}. ADR-027 dùng \texttt{web +
  muk\_web\_theme} (+ \texttt{ui\_core}); chỉ thêm \texttt{wujia\_core} khi mobile cần cấu hình
  Wujia thật.
\end{itemize}

\subsection*{Bảng BA lệch với code --- nhờ BA sửa lại}

\begin{center}
\footnotesize
\begin{tabularx}{\textwidth}{>{\raggedright\arraybackslash}p{4.6cm}>{\raggedright\arraybackslash}X}
\toprule
\textbf{BA ghi} & \textbf{Thực tế (code \texttt{12750a2} / UAT)} \\
\midrule
\texttt{wujia\_report} ``Đã xây dựng'' & Không tồn tại. Báo cáo nằm ở \texttt{wujia\_portal\_report}, \texttt{wujia.sale.supply.demand.report}, Dashboard Ninja. \\
\texttt{wujia\_delivery\_cost} ``Đã xây dựng'' & Không tồn tại. Cước nằm trong \texttt{wujia\_fleet} (bảng giá) và \texttt{wujia\_delivery}. \\
\texttt{wujia\_fleet} depend Odoo \texttt{fleet} & Không depend; \texttt{fleet} uninstalled trên UAT. Code depend \texttt{wujia\_core, wujia\_franchise, mail}. \\
\texttt{wujia\_delivery} depend \texttt{core, sale, franchise} & Code depend \texttt{wujia\_fleet, wujia\_sale, stock\_picking\_batch}. \\
\texttt{wujia\_sale} depend \texttt{core, franchise, sale\_management} & Code depend \texttt{sale, sale\_stock, stock, stock\_picking\_batch, wujia\_franchise}. \\
\texttt{portal\_base} depend \texttt{core, portal, website} & Code depend \texttt{wujia\_sale, wujia\_portal\_layout, portal, bus}. \texttt{portal\_layout} không có trong bảng. \\
\texttt{approval\_core} depend \texttt{wujia\_workflow\_core} & Module \texttt{wujia\_workflow\_core} không có dòng nào trong bảng. \\
(không có) & \texttt{wujia\_account}, \texttt{wujia\_franchise\_contract/\_inspection/\_operations}, 13 module \texttt{wujia\_portal\_*} khác, \texttt{portal\_order\_window}. \\
\bottomrule
\end{tabularx}
\end{center}

Theo quy tắc đã chốt từ Sprint 41, Dev \textbf{không ghi vào tài liệu BA}; bảng trên gửi BA tự sửa.

\section{Câu hỏi cần BA / chủ dự án chốt}

\begin{enumerate}
\item \textbf{Ai dùng backend trên điện thoại?} HQ duyệt, giám sát khu vực đi cửa hàng, kho, tài xế?
\item \textbf{3--5 màn ưu tiên đầu tiên} cho mobile core (kanban nào, form nào)? Câu trả lời
  quyết định phân hệ nào tách sớm.
\item \textbf{Có cần chọn nhiều bản ghi trên kanban mobile} (thao tác hàng loạt) không, hay chỉ
  cần chạm mở?
\item \textbf{Theme backend chốt là MuK?} \texttt{legion\_enterprise\_theme} đang nằm trong repo
  nhưng chưa cài. Nếu có ý định đổi thì phải chốt trước P1.
\item \textbf{Màn khảo sát \texttt{/franchise/inspection/do}}: dùng chủ yếu trên điện thoại? Nếu
  đúng, đây là ứng viên cầu nối mobile đầu tiên.
\item \textbf{Module \texttt{wujia\_portal\_<x>} còn được chứa phần duyệt/xử lý nghiệp vụ không?}
  Bảng BA đang xếp \texttt{portal\_return} như vậy; ADR-027 đề xuất đưa xuống \texttt{wujia\_return}.
  Phải chốt trước khi tách 7 phân hệ.
\item \textbf{BA sửa bảng ``0. Kiến trúc module''} theo mục đối chiếu \S\ref{sec:adr27-ba}.
\end{enumerate}

\section{Kết quả nghiệp vụ của phiên}

Chưa có thay đổi nào tới người dùng. Kết quả là \textbf{một bản đồ và một bộ luật}: biết chính xác
module nào đang nằm sai chỗ (7 module portal ôm nghiệp vụ), biết rằng nỗi lo ``mobile bị kéo lên
portal'' \emph{không} áp dụng cho franchise/sale/fleet/delivery/account, và có một lộ trình tách
từng bước với cách kiểm bằng số. Tiền lệ tách \texttt{wujia\_franchise} đã chạy trên UAT cho thấy
quy trình khả thi.
